\documentclass[12pt,letterpaper]{article}

% ---------------------------------------------------------------
% Paquetes
% ---------------------------------------------------------------
\usepackage[spanish,es-tabla]{babel}
\usepackage[utf8]{inputenc}
\usepackage[margin=2.1cm,top=1.6cm,bottom=1.4cm]{geometry}
\usepackage{xcolor}
\usepackage[most]{tcolorbox}
\usepackage{tikz}
\usepackage{titlesec}
\usepackage{enumitem}
\usepackage{parskip}
\usepackage[hidelinks]{hyperref}

\definecolor{sscobalt}{HTML}{1B6CA8}
\definecolor{sscoral}{HTML}{E2711D}
\definecolor{ssgray}{HTML}{4A4A4A}
\definecolor{ssgreen}{HTML}{2F9E44}
\definecolor{ssbg}{HTML}{F7F9FA}

\titleformat{\section}{\normalfont\Large\bfseries\color{sscobalt}}{\thesection.}{0.6em}{}

\setlength{\parindent}{0pt}
\setlength{\parskip}{3pt}

% ---------------------------------------------------------------
% Caja de segmento del guion: color, hora, título
% ---------------------------------------------------------------
\newtcolorbox{cuebox}[3]{
  enhanced,
  breakable,
  colback=white,
  colframe=#1,
  colbacktitle=#1,
  coltitle=white,
  fonttitle=\bfseries\sffamily,
  attach boxed title to top left={yshift=-2mm, xshift=3mm},
  boxed title style={boxrule=0pt, arc=2pt, top=2pt, bottom=2pt, left=7pt, right=7pt},
  title={\makebox[2.05cm][l]{\small #2}\,\normalsize #3},
  boxrule=1pt,
  arc=3pt,
  top=5.5mm, left=9pt, right=9pt, bottom=4pt,
  before skip=4pt, after skip=4pt,
}

% ---------------------------------------------------------------
% Etiqueta inline para acotaciones de escena
% ---------------------------------------------------------------
\newtcbox{\stagebox}{on line, colback=ssgray!12, colframe=ssgray!12,
  arc=2pt, boxrule=0pt, boxsep=1pt, left=5pt, right=5pt, top=2pt, bottom=2pt,
  fontupper=\small\itshape\color{ssgray!90!black}}

\newcommand{\stage}[1]{\stagebox{$\rightarrow$ #1}}

\pagestyle{empty}

\begin{document}

% ===================================================================
% PORTADA / ENCABEZADO
% ===================================================================
{\color{sscobalt}\rule{\linewidth}{2pt}}\\[0.35cm]
\begin{center}
  {\Huge\bfseries SteamSense Pro}\\[0.15cm]
  {\Large Guion de presentación --- pitch de negocio}
\end{center}
{\color{sscobalt}\rule{\linewidth}{1pt}}\\[0.4cm]

\begin{center}
\begin{tabular}{ccc}
\textbf{\textcolor{sscobalt}{\Large 5:00}} & \hspace{1.2cm} & \textbf{\textcolor{sscoral}{\Large 5}} \\
{\footnotesize duración total} & & {\footnotesize bloques} \\
\end{tabular}
\end{center}

\vspace{0.3cm}

% ---------------------------------------------------------------
% LÍNEA DE TIEMPO
% ---------------------------------------------------------------
\begin{center}
\begin{tikzpicture}[x=1cm, y=1cm]
  % barras proporcionales: 30s, 50s, 50s, 130s, 40s sobre 300s -> ancho total 15cm
  \definecolor{cA}{HTML}{1B6CA8}
  \definecolor{cB}{HTML}{E2711D}
  \definecolor{cC}{HTML}{2F9E44}
  \definecolor{cD}{HTML}{1B6CA8}
  \definecolor{cE}{HTML}{E2711D}

  \fill[cA] (0,0) rectangle (1.5,0.65);
  \fill[cB!85!white] (1.5,0) rectangle (4.0,0.65);
  \fill[cC!80!white] (4.0,0) rectangle (6.5,0.65);
  \fill[cD!60!white] (6.5,0) rectangle (13.0,0.65);
  \fill[cE!70!white] (13.0,0) rectangle (15.0,0.65);

  \draw[white, line width=1.2pt] (1.5,0) -- (1.5,0.65);
  \draw[white, line width=1.2pt] (4.0,0) -- (4.0,0.65);
  \draw[white, line width=1.2pt] (6.5,0) -- (6.5,0.65);
  \draw[white, line width=1.2pt] (13.0,0) -- (13.0,0.65);

  \node[white, font=\bfseries\footnotesize] at (0.75,0.325) {0:30};
  \node[white, font=\bfseries\footnotesize] at (2.75,0.325) {0:50};
  \node[white, font=\bfseries\footnotesize] at (5.25,0.325) {0:50};
  \node[white, font=\bfseries\footnotesize] at (9.75,0.325) {2:10};
  \node[black, font=\bfseries\footnotesize] at (14.0,0.325) {0:40};

  \node[font=\footnotesize, text=ssgray, anchor=north] at (0.75,-0.15) {Apertura};
  \node[font=\footnotesize, text=ssgray, anchor=north] at (2.75,-0.15) {Problema};
  \node[font=\footnotesize, text=ssgray, anchor=north] at (5.25,-0.15) {Estrategia};
  \node[font=\footnotesize, text=ssgray, anchor=north] at (9.75,-0.15) {3 casos reales};
  \node[font=\footnotesize, text=ssgray, anchor=north] at (14.0,-0.15) {Cierre};
\end{tikzpicture}
\end{center}

\vspace{0.2cm}

% ===================================================================
% BLOQUE 1 --- APERTURA
% ===================================================================
\begin{cuebox}{sscobalt}{0:00--0:30}{Apertura}
``Buenas tardes. Les voy a contar algo que seguramente les ha pasado: compran un juego, lo juegan un par de horas... y ahí se queda, sin terminar. Eso no es un problema del jugador --- es un problema de negocio. Esto es \textbf{SteamSense Pro}, y en los próximos minutos les voy a contar la historia de un jugador real, y cómo se la cambiamos \textit{antes} de que fuera demasiado tarde.''
\end{cuebox}

% ===================================================================
% BLOQUE 2 --- EL PROBLEMA
% ===================================================================
\begin{cuebox}{sscoral}{0:30--1:20}{El problema}
``Steam tiene miles de juegos de rol. Elegir cuál comprar es como entrar a una librería gigante sin saber qué te va a gustar --- y estos juegos no se leen en una tarde, te piden 40, 60, 100 horas de tu vida.

Y lo que vemos una y otra vez es siempre lo mismo: la gente compra... y no juega. Compra por una oferta, por una recomendación, por curiosidad --- y ese juego se queda ahí, sin abrir, acumulando polvo en la biblioteca.

Para el jugador, es plata y tiempo que no vuelven. Para el estudio que hizo ese juego, es un cliente que nunca vuelve por la secuela, ni por el DLC, ni lo recomienda a nadie.

Los dos pierden por la misma razón: \textbf{nadie se da cuenta a tiempo.}''
\end{cuebox}

% ===================================================================
% BLOQUE 3 --- LA ESTRATEGIA
% ===================================================================
\begin{cuebox}{ssgreen}{1:20--2:10}{La estrategia}
``Nuestra propuesta es simple: contarle a cada jugador su propia historia antes de que compre, para que no se arrepienta.

Antes de comprar, te mostramos lo que \textit{tú}, con tu forma de jugar, sí vas a terminar --- no lo más popular.

Después de comprar, si ya tienes el juego, te avisamos cuáles se te están yendo de las manos, a tiempo para retomarlos.

Y todo el tiempo, sin encuestas, sabemos qué tipo de jugador eres --- porque no todos compran, juegan, ni abandonan igual.

Esto es la diferencia entre un jugador que compra una vez y desaparece, y uno que confía en lo que le sugieres, porque acierta.''
\end{cuebox}

% ===================================================================
% BLOQUE 4 --- TRES CASOS REALES
% ===================================================================
\begin{cuebox}{sscobalt}{2:10--4:20}{Tres casos reales \normalfont\small(abre la app)}
``Se los muestro con tres casos reales. \stage{abrir la app}

\textbf{Primero, Vaps.} Tiene 84 juegos y más de 900 horas jugadas... pero solo 1 activo en las últimas dos semanas. Compró PAYDAY 2 y Gorky 17 y nunca los abrió. El sistema lo reconoce solo: es un \textbf{Coleccionista Masivo} --- compra mucho, pero su atención se reparte tan fino que casi nada llega a buen puerto. Por eso, en vez de ofrecerle otro juego de 80 horas, le recomendamos historias cortas, hechas para alguien que juega como él.

\stage{cambia de perfil en la app} \textbf{Ahora, juanmy03.} Apenas 3 juegos, ninguno de rol --- casi no hay rastro de qué le gusta. Y aquí el sistema es honesto: no inventa una recomendación perfecta solo por quedar bien. Le muestra apenas 30\% de compatibilidad, porque todavía no lo conoce lo suficiente. Eso también es valioso: le dice al estudio a quién le falta historial, antes de gastar en recomendarle algo a ciegas.

\stage{cambia de perfil en la app} \textbf{Y por último, Luistfmx10.} Miren esta cifra: 1{,}332 horas jugadas en total --- le ha dado la vida a Steam. Pero de esas horas, cero están en el único RPG que tiene, Black Desert, un juego de 94 horas de historia que nunca abrió. El sistema lo marca al instante como \textbf{Coleccionista Dormido}, con 98\% de probabilidad de perderlo para siempre. No es que no le guste jugar --- ahí están sus horas. Solo se alejó. Y ahora sabemos exactamente a quién reactivar, y con qué.

Tres jugadores completamente distintos. El mismo sistema, sin ayuda de nadie, entendió a los tres.''
\end{cuebox}

% ===================================================================
% BLOQUE 5 --- CIERRE
% ===================================================================
\begin{cuebox}{sscoral}{4:20--5:00}{Cierre}
``Esto no es una idea en papel. Ya existe, ya funciona, y se lo acabamos de mostrar en vivo, con jugadores reales.

Le sirve al jugador porque compra mejor y disfruta más lo que tiene. Y le sirve al estudio porque convierte algo invisible --- el abandono --- en algo que se puede ver, anticipar y remediar.

El siguiente paso es simple: \textbf{hacer esto con su propia comunidad de jugadores.} Gracias.''
\end{cuebox}

% ===================================================================
% TIPS
% ===================================================================
\begin{tcolorbox}[
  enhanced, breakable,
  colback=ssbg, colframe=ssgray!40,
  arc=3pt, boxrule=0.8pt,
  left=9pt, right=9pt, top=5pt, bottom=5pt,
  title={\bfseries\sffamily\color{ssgray} Consejos rápidos de ritmo},
  colbacktitle=ssbg, coltitle=ssgray, boxed title style={boxrule=0pt}
]
\begin{itemize}[leftmargin=1.1em, itemsep=0.12em, topsep=0.1em]
  \item El bloque de \textbf{tres casos reales} es, a propósito, el más largo --- ahí está tu venta real. No lo apures.
  \item Si vas corto de tiempo, el \textbf{cierre} es lo primero que puedes recortar a una sola frase.
  \item Practícalo una vez en voz alta con cronómetro --- a este ritmo de habla debería darte casi exactos los 5 minutos.
\end{itemize}
\end{tcolorbox}

\end{document}
